% \documentclass[11pt]{article}

% %%%% CUSTOM PREAMBLE
% \usepackage{./preamble}

% %% ADD MISC DEFINITIONS HERE %%%%%%%%%%%%%%%%%%%%%%%%%%%%%%%%%%%

% %%%%%%%%%%%%%%%%%%%%%%%%%%%%%%%%%%%%%%%%%%%%%%%%%%%%%
% %\addtolength{\topmargin}{-1cm}
% %\addtolength{\textheight}{2cm}

% \begin{document}

%\maketitle

\newpage
 %%*****************************************************
\noindent
\textbf{CS4268 AY2025-26 Sem 2}
\\
\noindent
Lecturer: Divesh Aggarwal
\hfill
Lecture 9                      %%% FILL IN LECTURE NUMBER HERE
\hfill
Date: March 19, 2026          %%% FILL IN LECTURE DATE HERE


\noindent
\rule{\textwidth}{1pt}

\medskip

%%%%%%%%%%%%%%%%%%%%%%%%%%%%%%%%%%%%%%%%%%%%%%%%%%%%%%%%%%%%%%%%
%% BODY OF SCRIBE NOTES GOES HERE
%%%%%%%%%%%%%%%%%%%%%%%%%%%%%%%%%%%%%%%%%%%%%%%%%%%%%%%%%%%%%%%%

%%%Write lecture topic above.
\textbf{Instructions}\\
\begin{itemize}
    \item Try to leave the preamble unchanged, only make changes to \texttt{main.tex} and  \texttt{references.bib}. Add images to the \texttt{images} folder. 
    \item Share one (zipped) directory preserving the directory structure on Canvas. 
    \item Prepare organized notes by elaborating on the broad items provided in different sections below based on the lecture material completing any missing details from lectures.
    \item Feel free to modify structure and/or order of contents as you see fit; the below is just a guideline.
\end{itemize} 



\section{Shor's Algorithm}

\kk{
Today is on Shor. The quantum subroutine is essentially simon over $Z_n$, a draft of which is given above if you need it. Please scribe in detail the number theoretic preliminaries (i believe there's a note somewhere by Aditya), the quantum part, and classical post-processing.
}

\subsection{Number-Theoretic Preliminaries}

We collect the number-theoretic background needed for Shor's algorithm. All integers are assumed to be positive unless stated otherwise. For $M \in \mathbb{Z}^+$ we write $\mathbb{Z}_M := \{1, 2, \dots, M-1\}$ and recall the standard fact: if $a \equiv a' \pmod{M}$ and $b \equiv b' \pmod{M}$, then $ab \equiv a'b' \pmod{M}$.

\subsubsection{Coprimality and Greatest Common Divisor}

\begin{definition}[Greatest Common Divisor]
Let $a$ and $b$ be integers, not both zero. The \emph{greatest common divisor} $\gcd(a,b)$ is the largest integer $d$ such that $d \mid a$ and $d \mid b$.
\end{definition}

\begin{definition}[Coprime]
Two integers $a$ and $b$ are \emph{coprime} (or \emph{relatively prime}) if $\gcd(a,b)=1$.
\end{definition}

\begin{example}
The divisors of $15$ are $\{1,3,5,15\}$ and the divisors of $28$ are $\{1,2,4,7,14,28\}$, so $\gcd(15,28)=1$ and they are coprime. In contrast, $\gcd(15,25)=5$, so $15$ and $25$ are not coprime.
\end{example}

\subsubsection{Bézout's Identity}

\begin{theorem}[Bézout's Identity]
\label{thm:bezout}
Let $a$ and $b$ be integers, not both zero. There exist integers $x$ and $y$ (the \emph{Bézout coefficients}) such that
\[
    ax + by = \gcd(a,b).
\]
\end{theorem}

For example, with $a=15$ and $b=28$ one can check that $x=15, y=-8$ gives $15 \cdot 15 + 28 \cdot (-8) = 225 - 224 = 1 = \gcd(15,28)$. Solutions are not unique: for any integer $k$, the pair $(x+kb,\, y-ka)$ is also a solution.

\begin{remark}
The equation $ax+by=c$ over the integers has integer solutions if and only if $\gcd(a,b) \mid c$. If $d = \gcd(a,b) \nmid c$ then $d$ divides the left-hand side but not the right, giving a contradiction.
\end{remark}

\subsubsection{The Extended Euclidean Algorithm}

The Extended Euclidean Algorithm (EEA) computes $\gcd(a,b)$ together with the Bézout coefficients. It augments the ordinary Euclidean algorithm with bookkeeping for the coefficients at each step.

\begin{algorithm}[H]
\caption{Extended Euclidean Algorithm}
\begin{algorithmic}[1]
\Require Integers $a \geq b \geq 0$
\Ensure Integers $g,x,y$ such that $g = \gcd(a,b)$ and $ax+by=g$
\State $r_0 \leftarrow a,\; r_1 \leftarrow b$
\State $s_0 \leftarrow 1,\; s_1 \leftarrow 0$
\State $t_0 \leftarrow 0,\; t_1 \leftarrow 1$
\State $i \leftarrow 1$
\While{$r_i > 0$}
    \State $q_i \leftarrow \lfloor r_{i-1}/r_i \rfloor$
    \State $r_{i+1} \leftarrow r_{i-1} - q_i r_i$
    \State $s_{i+1} \leftarrow s_{i-1} - q_i s_i$
    \State $t_{i+1} \leftarrow t_{i-1} - q_i t_i$
    \State $i \leftarrow i+1$
\EndWhile
\State \Return $r_{i-1},\; s_{i-1},\; t_{i-1}$
\end{algorithmic}
\end{algorithm}

The while loop runs $O(\log b)$ times and the total time complexity is $O(\log^2 a)$, dominated by integer multiplication.

\subsubsection{The Multiplicative Group \texorpdfstring{$\mathbb{Z}_M^*$}{Z*\_M}}

\begin{definition}[Commutative Group]
A \emph{commutative group} $(G, \cdot)$ is a set $G$ with a binary operation $\cdot$ satisfying closure, associativity, existence of an identity $e$, and existence of inverses. We restrict to finite commutative groups throughout.
\end{definition}

\begin{definition}[Multiplicative Group Modulo $M$]
For $M \in \mathbb{Z}^+$, define
\[
    \mathbb{Z}_M^* = \{a \in \mathbb{Z}_M \mid \gcd(a,M)=1\},
\]
equipped with multiplication modulo $M$.
\end{definition}

\begin{lemma}
$\mathbb{Z}_M^*$ is a group under multiplication modulo $M$.
\end{lemma}
\begin{proof}[Proof sketch]
Closure: $\gcd(a,M)=1$ and $\gcd(b,M)=1$ implies $\gcd(ab,M)=1$. Associativity follows from integer multiplication. The identity is $1$. For inverses: given $a \in \mathbb{Z}_M^*$, Bézout's identity (\cref{thm:bezout}) yields $ax + My = 1$ for some integers $x,y$; reducing modulo $M$ gives $ax \equiv 1 \pmod{M}$, so $x \bmod M$ is the multiplicative inverse of $a$.
\end{proof}

\begin{example}
$\mathbb{Z}_{15}^* = \{1,2,4,7,8,11,13,14\}$. Since $2 \times 8 = 16 \equiv 1 \pmod{15}$, the elements $2$ and $8$ are inverses of each other.
\end{example}

\subsubsection{Euler's Totient Function}

\begin{definition}[Euler's Totient Function]
For $M \in \mathbb{Z}^+$, Euler's totient $\varphi(M)$ counts the integers $1 \leq k \leq M$ coprime to $M$, i.e.\ $\varphi(M) = |\mathbb{Z}_M^*|$.
\end{definition}

\begin{theorem}
\label{thm:totient-prime}
If $M$ is prime, then $\varphi(M) = M-1$.
\end{theorem}
\begin{proof}
Every $k \in \mathbb{Z}_M$ satisfies $\gcd(k,M)=1$ since $M$ is prime, and $|\mathbb{Z}_M|=M-1$.
\end{proof}

\begin{theorem}
\label{thm:totient-pq}
If $p$ and $q$ are distinct primes and $M = pq$, then $\varphi(M) = (p-1)(q-1)$.
\end{theorem}
\begin{proof}
Among $\{1,\dots,pq\}$, the elements sharing a factor with $pq$ are the multiples of $p$ (there are $q$ of them) and the multiples of $q$ (there are $p$ of them), with $pq$ counted twice. So the count of elements not coprime to $pq$ is $p+q-1$, giving $\varphi(pq) = pq - (p+q-1) = (p-1)(q-1)$.
\end{proof}

\subsubsection{Euler's Theorem}

\begin{theorem}[Euler's Theorem]
\label{thm:euler}
If $\gcd(a,M)=1$, then $a^{\varphi(M)} \equiv 1 \pmod{M}$.
\end{theorem}
\begin{proof}
Let $\mathbb{Z}_M^* = \{x_1,\dots,x_k\}$ with $k=\varphi(M)$. Since $a \in \mathbb{Z}_M^*$, multiplication by $a$ permutes the elements. Thus
\[
    \prod_{i=1}^k x_i \equiv \prod_{i=1}^k (ax_i) \equiv a^k \prod_{i=1}^k x_i \pmod{M}.
\]
Since each $x_i$ is coprime to $M$, their product is invertible modulo $M$; cancelling it gives $a^{\varphi(M)} \equiv 1 \pmod{M}$.
\end{proof}

\subsubsection{Order and Generators}

\begin{definition}[Order of an Element]
\label{def:order}
For $a \in \mathbb{Z}_M^*$, the \emph{order} $\mathrm{ord}_M(a)$ is the smallest positive integer $r$ such that $a^r \equiv 1 \pmod{M}$.
\end{definition}

\begin{lemma}[Order Divides Totient]
\label{lem:order-divides}
For any $a \in \mathbb{Z}_M^*$, $\mathrm{ord}_M(a)$ divides $\varphi(M)$.
\end{lemma}
\begin{proof}
Let $r = \mathrm{ord}_M(a)$. Write $\varphi(M) = qr + s$ with $0 \leq s < r$. By Euler's Theorem, $1 \equiv a^{\varphi(M)} = (a^r)^q \cdot a^s \equiv a^s \pmod{M}$. Since $r$ is minimal and $0 \leq s < r$, we must have $s=0$, so $r \mid \varphi(M)$.
\end{proof}

\begin{corollary}
\label{cor:order-prime}
Let $p$ be prime. Then $\mathrm{ord}_p(a)$ divides $p-1$ for all $a \in \mathbb{Z}_p^*$.
\end{corollary}

\begin{definition}[Cyclic Group and Generator]
A commutative group $(G,\cdot)$ is \emph{cyclic} if there exists a \emph{generator} $g \in G$ such that every element of $G$ is a power of $g$. We write $G = \langle g \rangle$.
\end{definition}

In $\mathbb{Z}_M^*$, a generator is called a \emph{primitive root modulo $M$}. For example, $2$ generates $\mathbb{Z}_5^* = \{1,2,3,4\}$ since $2^1=2$, $2^2=4$, $2^3 \equiv 3$, $2^4 \equiv 1 \pmod{5}$.

\begin{lemma}
\label{lem:totient-sum}
For any $M \geq 1$, $\displaystyle\sum_{d \mid M} \varphi(d) = M$.
\end{lemma}
\begin{proof}
Consider the $M$ fractions $\frac{1}{M},\frac{2}{M},\dots,\frac{M}{M}$. Each reduces to an irreducible fraction $\frac{k}{d}$ with $d \mid M$. For a fixed divisor $d$, the number of fractions reducing to denominator $d$ is exactly $\varphi(d)$. Summing over all divisors gives $M$.
\end{proof}

\begin{lemma}
\label{lem:Np}
Let $p$ be prime and $d \mid (p-1)$. Let $N_p(d)$ denote the number of elements in $\mathbb{Z}_p^*$ of order exactly $d$. Then $N_p(d) \in \{0, \varphi(d)\}$.
\end{lemma}
\begin{proof}
If $N_p(d)>0$, let $a \in \mathbb{Z}_p^*$ have order $d$. The elements $1, a, a^2, \dots, a^{d-1}$ are $d$ distinct roots of $x^d \equiv 1 \pmod{p}$. Since a degree-$d$ polynomial over a prime field has at most $d$ roots, these are all of them. Any element of order $d$ must be one of these, and $a^k$ has order $d$ iff $\gcd(k,d)=1$, of which there are $\varphi(d)$ choices.
\end{proof}

\begin{theorem}[Existence of Generators]
\label{thm:generators}
For any prime $p$, the multiplicative group $\mathbb{Z}_p^*$ is cyclic.
\end{theorem}
\begin{proof}
By \cref{lem:Np}, $N_p(d) \leq \varphi(d)$ for all $d \mid (p-1)$. Every element of $\mathbb{Z}_p^*$ has some order dividing $p-1$ (\cref{lem:order-divides}), so $\sum_{d \mid (p-1)} N_p(d) = p-1$. By \cref{lem:totient-sum}, $\sum_{d \mid (p-1)} \varphi(d) = p-1$ as well. Since each term satisfies $N_p(d) \leq \varphi(d)$ and the sums are equal, we must have $N_p(d) = \varphi(d)$ for every divisor $d$. In particular $N_p(p-1) = \varphi(p-1) \geq 1$, so there exists an element of order $p-1$, i.e.\ a generator.
\end{proof}















% \bibliographystyle{alpha}
% \bibliography{main.bib}
% \end{document}
